\section{The Hybrid $S_N$-Diffusion Method} \label{sec:hybrid}

The hybrid method is an iterative method that applies the $S_N$ discrete ordinates neutron
transport method on subdomains containing the control rod to generate pointwise transport
corrections for the neutron diffusion method. It consists of an $S_N$ neutron transport solver and
a neutron diffusion solver iteratively coupled through transport correction and boundary source
terms. This section presents mathematical formulations of the
\gls{SAAF} $S_N$ equations and the drift transport correction term for the
neutron diffusion equations, and the computational algorithm for the hybrid method.

\subsection{Self-Adjoint Angular Flux (SAAF) $S_N$ Method} \label{sec:saaf}

For the $S_N$ method implementation in this work, we implemented a second-order linear \gls{SAAF}
formulation with void stabilization treatment developed by Wang et al.\ \cite{wangDiffusionAccelerationSchemes2014}.
This formulation is compatible with the highly efficient and scalable HYPRE-BoomerAMG
preconditioner \cite{hypre_hypre_2022} available through \gls{MOOSE}. Implementing the $S_N$ method
in Moltres enables straightforward interfacing with the existing neutron diffusion solver for
nonlinear diffusion acceleration and for the subsequent hybrid method implementation.

The time-independent multigroup neutron transport equation with boundary sources for the angular
neutron flux $\Psi_g$ defined on the 3-D spatial domain $\mathcal{D}$ and 2-D unit sphere angular
domain $\mathcal{S}$ with $G$ neutron energy groups is
%
\begin{align}
  \hat{\Omega}\cdot\nabla\Psi_g(\vec{r},\hat{\Omega}) + \Sigma_{t,g}(\vec{r})
  \Psi_g(\vec{r},\hat{\Omega})
  =& \sum^G_{g'=1}\int_\mathcal{S} \Sigma_s^{g'\rightarrow g}(\vec{r}, \hat{\Omega}'\rightarrow\hat{\Omega})
  \Psi_{g'}(\vec{r},\hat{\Omega}')d\hat{\Omega}' \nonumber \\
   &+ \frac{1}{4\pi}\chi_{p,g}(1-\beta)\sum^G_{g'=1} \nu\Sigma_{f,g'}(\vec{r}) \phi_{g'}(\vec{r})
  + \frac{1}{4\pi}\sum^I_{i=1}\chi_{d,g}\lambda_i C_i(\vec{r}) \label{eq:mg-nte} \\
  \Psi_g(\vec{r},\hat{\Omega}) =& \Psi^\text{inc}_g(\vec{r},\hat{\Omega}) +
  \alpha^s_g\Psi_g(\vec{r},\hat{\Omega}_r)
  \mbox{ for } \vec{r} \in \partial\mathcal{D} \mbox{ and } \hat{\Omega}\cdot\hat{n}_b < 0, \label{eq:bc}
  \shortintertext{where}
  \vec{r} =& \mbox{ spatial coordinates,} \nonumber \\
  \hat{\Omega} =& \mbox{ direction of neutron travel,} \nonumber \\
  \Sigma_{t,g}(\vec{r}) =& \mbox{ macroscopic total cross section in group $g$,} \nonumber \\
  \Sigma_s^{g'\rightarrow g}(\vec{r},\hat{\Omega}'\rightarrow \hat{\Omega}) =&\nonumber
  \mbox{ macroscopic scattering cross section in group $g$,} \nonumber \\
  \chi_{p,g} =& \mbox{ fission neutron spectrum for group $g$,} \nonumber \\
  \nu =& \mbox{ number of neutrons produced per fission reaction,} \nonumber \\
  \Sigma_{f,g}(\vec{r}) =& \mbox{ macroscopic fission cross section in group $g$,} \nonumber \\
  \chi_{p,g}(\vec{r}) &= \mbox{prompt fission neutron spectrum in group $g$,} \nonumber\\
  \beta(\vec{r},) &= \sum^I_{i=1} \beta_i(\vec{r},t) = \mbox{total delayed neutron fraction,} \nonumber\\
  \chi_{d,g}(\vec{r}) &= \mbox{delayed fission neutron spectrum in group $g$,} \nonumber\\
  \lambda_i(\vec{r}) &= \mbox{decay constant of precursor group $i$,} \nonumber\\
  C_i(\vec{r}) &= \mbox{delayed neutron precursor concentration for group $i$,} \nonumber\\
  \Psi^\text{inc}_g(\vec{r}) &= \mbox{incident surface source in group $g$,} \nonumber\\
  \alpha^s_g &= \mbox{specular reflectivity on }\partial \mathcal{D} \mbox{ for group }g, \nonumber\\
  \hat{\Omega}_r &= \hat{\Omega}-2(\hat{\Omega}\cdot \hat{n}_b)\hat{n}_b, \nonumber\\
  \hat{n}_b &= \mbox{outward unit normal vector on the boundary.}\nonumber
\end{align}

We define the following vector forms:
%
\begin{gather}
  \bm{\Psi} \equiv
  \begin{bmatrix}
    \Psi_1 \\
    \Psi_2 \\
    \vdots \\
    \Psi_G
  \end{bmatrix},
  \bm{\Phi} \equiv \int_S \bm{\Psi}d\hat{\Omega} \equiv
  \begin{bmatrix}
    \phi_1 \\
    \phi_2 \\
    \vdots \\
    \phi_G
  \end{bmatrix},
  \bm{C} \equiv
  \begin{bmatrix}
    C_1 \\
    C_2 \\
    \vdots \\
    C_G
  \end{bmatrix},
\end{gather}
%
and the following operators:
%
\begin{gather}
  \mathbb{L}_1\bm{\Psi} \equiv
  \begin{bmatrix}
    \hat{\Omega}\cdot\nabla\Psi_1 \\
    \hat{\Omega}\cdot\nabla\Psi_2 \\
    \vdots \\
    \hat{\Omega}\cdot\nabla\Psi_G \\
  \end{bmatrix},
  \mathbb{L}_2\bm{\Psi} \equiv
  \begin{bmatrix}
    \Sigma_{t,1}\Psi_1 \\
    \Sigma_{t,2}\Psi_2 \\
    \vdots \\
    \Sigma_{t,G}\Psi_G
  \end{bmatrix},
  \mathbb{L}\bm{\Psi} \equiv \mathbb{L}_1\bm{\Psi} + \mathbb{L}_2\bm{\Psi},
  \mathbb{L}^*\bm{\Psi} \equiv -\mathbb{L}_1\bm{\Psi} + \mathbb{L}_2\bm{\Psi}, \nonumber \\
  \mathbb{S}\bm{\Psi} \equiv
  \begin{bmatrix}
    \sum^G_{g'=1}\int_S \Sigma_s^{g'\rightarrow 1}\Psi_{g'}d\hat{\Omega} \\
    \sum^G_{g'=2}\int_S \Sigma_s^{g'\rightarrow 2}\Psi_{g'}d\hat{\Omega} \\
    \vdots \\
    \sum^G_{g'=G}\int_S \Sigma_s^{g'\rightarrow G}\Psi_{g'}d\hat{\Omega}
  \end{bmatrix},
  \mathbb{B}\bm{\Psi} \equiv
  \begin{bmatrix}
    \alpha^s_1\Psi_1(\hat{\Omega}_r) \\
    \alpha^s_2\Psi_2(\hat{\Omega}_r) \\
    \vdots \\
    \alpha^s_G\Psi_G(\hat{\Omega}_r)
  \end{bmatrix}, \nonumber \\
  \mathbb{F}\bm{\Psi} \equiv \frac{1}{4\pi}\mathbb{F}_0\bm{\Psi} \equiv
  \begin{bmatrix}
    \frac{1}{4\pi}\chi_{p,1}(1-\beta)\sum^G_{g'=1}\nu\Sigma_{f,g'}\phi_{g'} \\
    \frac{1}{4\pi}\chi_{p,2}(1-\beta)\sum^G_{g'=1}\nu\Sigma_{f,g'}\phi_{g'} \\
    \vdots \\
    \frac{1}{4\pi}\chi_{p,G}(1-\beta)\sum^G_{g'=1}\nu\Sigma_{f,g'}\phi_{g'}
  \end{bmatrix},
  \mathbb{C}\bm{C} \equiv \frac{1}{4\pi}\mathbb{C}_0\bm{C} \equiv
  \begin{bmatrix}
    \frac{1}{4\pi}\sum^I_{i=1}\chi_{d,1} \lambda_i C_i \\
    \frac{1}{4\pi}\sum^I_{i=1}\chi_{d,2} \lambda_i C_i \\
    \vdots \\
    \frac{1}{4\pi}\sum^I_{i=1}\chi_{d,G} \lambda_i C_i
  \end{bmatrix}, \nonumber
\end{gather}
%
to reexpress Eq.\ \ref{eq:mg-nte} and \ref{eq:bc} as
%
\begin{gather}
  \mathbb{L}\bm{\Psi} = \mathbb{S}\bm{\Psi} + \bm{Q}_f, \label{eq:nte-vec} \\
  \bm{\Psi} = \bm{\Psi}^\text{inc} + \mathbb{B}\bm{\Psi}.
  \shortintertext{where}
  \bm{Q}_f = \mathbb{F}\bm{\Psi} + \mathbb{C}\bm{C},
\end{gather}

We also introduce inner product notations
to express Eq.\ \ref{eq:nte-vec} and the eventual \gls{SAAF} formulation in their weak forms and
to evaluate residuals for the \gls{FEM}. We define the inner product consisting of volume
integrals on $\mathcal{D}\otimes\mathcal{S}$ as
%
\begin{gather}
  (\bm{a}, \bm{b})_{\mathcal{D}\otimes\mathcal{S}} \equiv
  \sum^G_{g=1}\int_\mathcal{S}d\hat{\Omega}\int_\mathcal{D}d\vec{r}\
  a_g(\vec{r},\hat{\Omega}) b_g(\vec{r},\hat{\Omega}), \label{eq:weak-domain}
\end{gather}
%
where $\bm{a}$ and $\bm{b}$ represent multigroup function vectors such as $\bm{\Psi}$. We also
define boundary integrals on $\partial\mathcal{D}\otimes\mathcal{S}$ as
%
\begin{gather}
  \langle\bm{a},\bm{b}\rangle^\pm_{\partial\mathcal{D}\otimes\mathcal{S}} \equiv
  \sum^G_{g=1}\int_{\partial\mathcal{D}}d\vec{r}
  \int_{\mathcal{S}^{\pm}}d\hat{\Omega}\ |\hat{\Omega}\cdot\hat{n}_b|
  a_g(\vec{r},\hat{\Omega}) b_g(\vec{r},\hat{\Omega}), \label{eq:weak-boundary}
\end{gather}
%
where the $\pm$ sign depends on the signedness of $\hat{\Omega}\cdot\hat{n}_b$ at the boundary.
%
In accordance with standard weak form derivation procedure, Eq.\ \ref{eq:nte-vec} is multiplied
by a test function $\bm{\Psi}^*$ and integrate throughout over the $\mathcal{D}$ and $\mathcal{S}$
domains
%
\begin{gather}
  \left(\bm{\Psi}^*,\mathbb{L}\bm{\Psi}\right)_{\mathcal{D}\otimes\mathcal{S}} =
  \left(\bm{\Psi}^*,\mathbb{S}\bm{\Psi}\right)_{\mathcal{D}\otimes\mathcal{S}} +
  \left(\bm{\Psi}^*,\bm{Q}\right)_{\mathcal{D}\otimes\mathcal{S}}.
\end{gather}
%
Applying integration by parts on the streaming term,
$\left(\bm{\Psi}^*,\mathbb{L}\bm{\Psi}\right)+\langle\bm{\Psi}^*,\bm{\Psi}\rangle^- =
\left(\mathbb{L}^*\bm{\Psi}^*,\bm{\Psi}\right)+\langle\bm{\Psi}^*,\bm{\Psi}\rangle^+$, gives
%
\begin{gather}
  \left(\mathbb{L}^*\bm{\Psi}^*,\bm{\Psi}\right)_{\mathcal{D}\otimes\mathcal{S}} +
  \langle\bm{\Psi}^*,\bm{\Psi}\rangle_{\partial\mathcal{D}\otimes\mathcal{S}}^+ -
  \langle\bm{\Psi}^*,\bm{\Psi}\rangle_{\partial\mathcal{D}\otimes\mathcal{S}}^- =
  \left(\bm{\Psi}^*,\mathbb{S}\bm{\Psi}\right)_{\mathcal{D}\otimes\mathcal{S}} +
  \left(\bm{\Psi}^*,\bm{Q}\right)_{\mathcal{D}\otimes\mathcal{S}}. \label{eq:nte-weak}
\end{gather}

The streamline-upwind Petrov-Galerkin stabilization technique \cite{brooks_streamline_1982}, first
applied to \gls{SAAF} $S_N$ equations by Wang et al.\ \cite{wangDiffusionAccelerationSchemes2014}, provides
numerical stability in near-void regions and is given as
%
\begin{multline}
  \left(\mathbb{L}_1\bm{\Psi}^*,
  \left(\bm{\tau}\mathbb{L}_1-\mathbb{I}+\bm{\tau}\mathbb{L}_2\right)\bm{\Psi}\right)_{\mathcal{D}\otimes\mathcal{S}} +
  \langle\bm{\Psi}^*,\bm{\Psi}\rangle_{\partial\mathcal{D}\otimes\mathcal{S}}^+ -
  \langle\bm{\Psi}^*,\bm{\Psi}\rangle_{\partial\mathcal{D}\otimes\mathcal{S}}^- +
  \left(\mathbb{L}_2\bm{\Psi}^*,\bm{\Psi}\right)_{\mathcal{D}\otimes\mathcal{S}} \\
  = \left(\left(\mathbb{I}+\bm{\tau}\mathbb{L}_1\right)\bm{\Psi}^*,\mathbb{S}\bm{\Psi}\right)_{\mathcal{D}\otimes\mathcal{S}} +
  \left(\left(\mathbb{I}+\bm{\tau}\mathbb{L}_1\right)\bm{\Psi}^*,\bm{Q}\right)_{\mathcal{D}\otimes\mathcal{S}}
  \label{eq:saaf-vt}
\end{multline}
%
\begin{gather}
  \shortintertext{where}
  \begin{align*}
    \bm{\tau} &\equiv
    \begin{bmatrix}
      \tau_1 \\
      \tau_2 \\
      \vdots \\
      \tau_G
    \end{bmatrix} \\
    \tau_g &=
    \begin{cases}
      \frac{1}{c\Sigma_{t,g}} \mbox{ for } ch\Sigma_{t,g} \geq \varsigma \\
      \frac{h}{\varsigma} \mbox{ for } ch\Sigma_{t,g} < \varsigma
    \end{cases}, & \\
    h &= \mbox{mesh element size,} & \\
    c &= \mbox{maximum stabilization factor,} & \\
    \varsigma &= \mbox{void constant.} &
  \end{align*}
\end{gather}

In \gls{FEM}, the spatial domain is discretized into finite, non-overlapping mesh
elements $\mathcal{D}_e$ such that $\mathcal{D} = \cup_{e\in\mathbb{T}_h}\mathcal{D}_e$, where
$\mathbb{T}_h$ denotes the index set of elements discretizing $\mathcal{D}$ and $\mathcal{D}_e$
is the subdomain covered by the element of index $e$. Similarly, the outermost boundary
$\partial\mathcal{D}$ is discretized by non-overlapping mesh sides indexed by $s$. Within each mesh
element, any spatially varying function can be approximated with basis functions
%
\begin{gather}
  f(\vec{r}) = \sum^N_{i=1}f_i b_i(\vec{r}),
  \shortintertext{where}
  \begin{align*}
    N &= \mbox{degrees of freedom in the mesh element,} \\
    b_i &= \mbox{basis function,} \\
    f_i &= \mbox{coefficient corresponding to $b_i$ in approximating $f$}.
  \end{align*}
\end{gather}

Thus, we can define the inner products
%
\begin{gather}
  \left(\bm{a},\bm{b}\right)_\mathcal{D} \equiv \int_\mathcal{D}d\vec{r}
  \sum_{i\in E_e}a_{g,i}b_i(\vec{r})\sum_{j\in E_e}b_{g,j}b_j(\vec{r}), \\
  \left(\bm{a},\bm{b}\right)_{\partial\mathcal{D}} \equiv
  \sum_{s\in\partial\mathcal{D}}\int_s d\vec{r}\sum_{i\in E_s}a_{g,i}b_i(\vec{r})\sum_{j\in E_s}
  b_{g,j}b_j(\vec{r}),
  \shortintertext{where}
  \begin{align*}
    E_e &= \mbox{index set of the local basis functions in element $e$,} \\
    E_s &= \mbox{index set of the local basis functions in element side $s$,}
  \end{align*}
\end{gather}
%
to apply finite element spatial discretization.

The $S_N$ method discretizes the angular dependence using a collocation method along $N_d$ discrete
directions $\hat{\Omega}_d$ with weights $w_d$. For this work, we chose the commonly used
level-symmetric quadrature set \cite{wang_rattlesnake_2018} for ease of implementation,
particularly when applying reflective boundary conditions on planes normal to one of the three
Cartesian axes. Thus we approximate the zeroth and higher-order moments of neutron flux using real
spherical harmonic functions $Y_{l,m}$ of degree $l$ and order $m$ as follows
%
\begin{gather}
  \phi_{g,0,0}(\vec{r}) = \phi_g(\vec{r}) = \int_\mathcal{S} \Psi_g(\vec{r}) d\hat{\Omega} \approx \sum^{N_d}_{d=1}w_d
  \Psi_{g,d}(\vec{r}), \label{eq:0th-mmt} \\
  \phi_{g,l,m}(\vec{r}) = \int_\mathcal{S} Y_{l,m}(\hat{\Omega})\Psi_g(\vec{r},\hat{\Omega})d\hat{\Omega} \approx
  \sum^{N_d}_{d=1}w_d Y_{l,m}(\hat{\Omega}_d)\Psi_{g}(\vec{r},\hat{\Omega}_d) \label{eq:nth-mmt}
  \shortintertext{where}
  \begin{align*}
    \Psi_{g,d}(\vec{r}) =& \Psi_g(\vec{r},\hat{\Omega}_d), \\
    w =& \sum^{N_d}_{d=1} w_d\ (= 8 \mbox{ for 3-D level-symmetric set}). \\
  \end{align*}
\end{gather}
%

The weak form of the fully discretized multigroup \gls{SAAF} $S_N$ equations with void
stabilization treatment is
%
\begin{multline}
  \sum^{N_d}_{d=1}w_d\left(\hat{\Omega}_d\cdot\nabla\Psi^*_{g,d},\tau_g\hat{\Omega}
  \cdot\nabla\Psi_{g,d}-(1-\tau_g\Sigma_{t,g})\Psi_{g,d}\right)_\mathcal{D}
  + \sum^{N_d}_{d=1}w_d\left(\Psi^*_{g,d},\Sigma_{t,g}\Psi_{g,d}\right)_\mathcal{D} \\
  = \sum^{N_d}_{d=1}w_d\left(\Psi^*_{g,d}+\tau_g\hat{\Omega}_d\cdot\nabla\Psi^*_{g,d},
  \sum^G_{g'=1}\sum^L_{l=0}\Sigma^{g'\rightarrow g}_{s,l}\sum^l_{m=-l}
  \frac{2l+1}{w}Y_{l,m}(\hat{\Omega}_d)\phi_{g',l,m}\right)_\mathcal{D} \\
  + \sum^{N_d}_{d=1}w_d\left(\Psi^*_{g,d}+\tau_g\hat{\Omega}_d\cdot\nabla\Psi^*_{g,d},
  \frac{1}{w}\bar{\chi}_g\sum^G_{g'=1}\nu\Sigma_{f,g'}\phi_{g'}\right)_\mathcal{D} \\
  + \sum^{N_d}_{d=1}w_d\left(\Psi^*_{g,d}+\tau_g\hat{\Omega}_d\cdot\nabla\Psi^*_{g,d},
  \frac{1}{w}\sum ^I_{i=1}\chi_{d,g}\lambda_i C_i\right)_\mathcal{D}
  \label{eq:weak-form}
\end{multline}
%
where $\bar{\chi}_g=\chi_{p,g}\left(1-\beta\right)$ for problems with \glspl{DNP}
and $\bar{\chi}_g=\chi_{g}$ for problems without \glspl{DNP}. $\vec{r}$ dependence is dropped for
brevity.
Both prompt fission and delayed neutron source terms scale by the inverse of the
multiplication factor $\frac{1}{k}$ for $k$-eigenvalue problems.

The weak forms of the vacuum boundary, boundary source, and reflecting boundary terms 
%
\begin{gather}
  \mbox{Vacuum } =
  \begin{cases}
    \sum^G_{g=1}\sum^{N_d}_{d=1}w_d\left(\Psi^*_{g,d},
    \hat{\Omega}_d\cdot\hat{n}_b\Psi_{g,d}\right)_{\partial\mathcal{D}}
    & \forall\hat{\Omega}\cdot\hat{n}_b>0,\vec{r}\in\partial\mathcal{D} \\
    0
    & \forall\hat{\Omega}\cdot\hat{n}_b<0,\vec{r}\in\partial\mathcal{D}
  \end{cases},
\end{gather}
%
\begin{gather}
  \vcenter{\hbox{\shortstack{Boundary \\ source }}} =
  \begin{cases}
    \sum^G_{g=1}\sum^{N_d}_{d=1}w_d\left(\Psi^*_{g,d},
    \hat{\Omega}_d\cdot\hat{n}_b\Psi_{g,d}\right)_{\partial\mathcal{D}}
    & \forall\hat{\Omega}\cdot\hat{n}_b>0,\vec{r}\in\partial\mathcal{D} \\
    \sum^G_{g=1}\sum^{N_d}_{d=1}w_d\left(\Psi^*_{g,d},
    \hat{\Omega}_d\cdot\hat{n}_b\Psi^\text{inc}_{g,d}\right)_{\partial\mathcal{D}}
    & \forall\hat{\Omega}\cdot\hat{n}_b<0,\vec{r}\in\partial\mathcal{D}
  \end{cases}, \label{eq:boundary-source}
\end{gather}
%
\begin{gather}
  \mbox{Reflecting } =
  \begin{cases}
    \sum^G_{g=1}\sum^{N_d}_{d=1}w_d\left(\Psi^*_{g,d},
    \hat{\Omega}_d\cdot\hat{n}_b\Psi_{g,d}\right)_{\partial\mathcal{D}}
    & \forall\hat{\Omega}\cdot\hat{n}_b>0,\vec{r}\in\partial\mathcal{D}_s \\
    \sum^G_{g=1}\sum^{N_d}_{d=1}w_d\left(\Psi^*_{g,d},
    \hat{\Omega}_d\cdot\hat{n}_b\Psi_{g,d_r}\right)_{\partial\mathcal{D}}
    & \forall\hat{\Omega}\cdot\hat{n}_b<0,\vec{r}\in\partial\mathcal{D}_s
  \end{cases}, \label{eq:reflecting-bc}
  \shortintertext{where}
  \hat{\Omega}_{d_r} = \hat{\Omega}_d - 2(\hat{\Omega}_d\cdot\hat{n}_b)\hat{n}_b = \mbox{ direction
  of incident neutron travel}. \nonumber
\end{gather}

\subsection{Transport Correction Formulation} \label{sec:transport-correction}

The conventional approach for determining diffusion coefficients for the neutron diffusion
equations in each subregion involves
running a high-fidelity neutron transport simulation to tally region-wide estimates of the neutron
transport cross section. The transport cross section formulation derives from the $P_1$
approximation of the neutron transport equation with isotropic sources \cite{bellNuclearReactorTheory1970a}.
Each defined subregion has a constant diffusion coefficient value. However, the neutron diffusion
equation is only valid in regions of high scattering-to-removal ratios with, at most, linearly
anisotropic scattering and small flux gradients. These conditions do not hold near or within
control rods, near interfaces of neighboring materials with highly dissimilar neutronic properties,
and in materials with significant scattering contributions from light nuclei.

Transport corrections to allow neutron diffusion methods to reproduce high-fidelity flux solutions
can be introduced by modifying the neutron diffusion equations. This work explores the use of
drift terms to apply transport corrections to the neutron diffusion
equations.

Drift correction terms may be added to neutron diffusion or equivalent low-order equations to
correct them such that the modified equations can reproduce the transport solution upon reaching
iterative convergence. Drift terms are prominent in existing literature as transport correction
terms for \gls{NDA} schemes \cite{adams_fast_2002, wangDiffusionAccelerationSchemes2014,
sun_discrete-ordinates_2024} or nodal diffusion methods
\cite{smith_nodal_1983, smith_assembly_1986}. \gls{NDA} schemes fall under \gls{HOLO} methods
\cite{chacon_multiscale_2017} similar to the \gls{QD} method
for accelerating a neutron transport
method (high-order) with a modified neutron diffusion method (low-order). Nodal
diffusion methods apply drift terms to rectify the incoming and outgoing neutron
currents between adjacent nodes.

The general form of a drift correction term for the neutron diffusion equations is a
first-order derivative $\nabla\cdot \vec{D}_g\phi_g$ (note the vector notation).
The exact forms depend on the discretization schemes of the high- and low-order
equations. For this work, we adopt drift terms derived by Wang et al.\
\cite{wangDiffusionAccelerationSchemes2014, wang_rattlesnake_2018} originally developed for \gls{NDA} of the
\gls{SAAF} $S_N$ equations with void treatment. We start by integrating Eq.\ \ref{eq:saaf-vt}
over the 2-D unit sphere angular domain $\mathcal{S}$ and replacing $\bm{\Psi}^*$ with isotropic
test functions $\bm{\Phi}^*$ to obtain the neutron balance equation
%
\begin{multline}
%  \left(\bm{\Phi}^*,\frac{\partial}{\partial t}\left(\frac{\bm{\Phi}}{\bm{v}}\right)\right)_\mathcal{D}
%  +
  \left(\mathbb{L}_2\bm{\Phi}^*,\bm{\Phi}\right)_\mathcal{D}
  - \left(\nabla\bm{\Phi}^*,\vec{\bm{J}}\right)_\mathcal{D}
  + \left(\bm{\Phi}^*,\bm{J}^\text{out}\right)_{\partial\mathcal{D}_v} + \\
  \left(\bm{\tau}\nabla\bm{\Phi}^*,
    \nabla\cdot(\vec{\vec{\bm{E}}}\bm{\Phi}) +\mathbb{L}_2\vec{\bm{J}} -
    \mathbb{S}_1\vec{\bm{J}}\right)_\mathcal{D}
  = \left(\bm{\Phi}^*,\mathbb{S}_0\bm{\Phi}\right)_\mathcal{D}
  + \left(\bm{\Phi}^*,\bm{Q}_0\right)_\mathcal{D},
\end{multline}
%
\begin{gather}
  \shortintertext{where}
  \vec{\bm{J}} \equiv
  \begin{bmatrix}
    \vec{J}_1 \\
    \vec{J}_2 \\
    \vdots \\
    \vec{J}_G
  \end{bmatrix},
  \bm{J}^\text{out} \equiv
  \begin{bmatrix}
    \int_{|\hat{\Omega}\cdot\hat{n}_b|>0}|\hat{\Omega}\cdot\hat{n}_b|\Psi_1 d\hat{\Omega} \\
    \int_{|\hat{\Omega}\cdot\hat{n}_b|>0}|\hat{\Omega}\cdot\hat{n}_b|\Psi_2 d\hat{\Omega} \\
    \vdots \\
    \int_{|\hat{\Omega}\cdot\hat{n}_b|>0}|\hat{\Omega}\cdot\hat{n}_b|\Psi_G d\hat{\Omega}
  \end{bmatrix},
  \bm{J}^\text{inc} \equiv
  \begin{bmatrix}
    J^\text{inc}_1 \\
    J^\text{inc}_2 \\
    \vdots \\
    J^\text{inc}_G
  \end{bmatrix},
  \vec{\vec{\bm{E}}} \equiv
  \begin{bmatrix}
    \frac{\int_\mathcal{S} \hat{\Omega}\hat{\Omega}\Psi_1 d\hat{\Omega}}{\int_\mathcal{S}\Psi_1 d\hat{\Omega}} \\
    \frac{\int_\mathcal{S} \hat{\Omega}\hat{\Omega}\Psi_2 d\hat{\Omega}}{\int_\mathcal{S}\Psi_2 d\hat{\Omega}} \\
    \vdots \\
    \frac{\int_\mathcal{S} \hat{\Omega}\hat{\Omega}\Psi_G d\hat{\Omega}}{\int_\mathcal{S}\Psi_G d\hat{\Omega}}
  \end{bmatrix}, \nonumber \\
  \mathbb{S}_0 \bm{\Phi} \equiv
  \begin{bmatrix}
    \sum^G_{g'=1}\Sigma_s^{g'\rightarrow 1}\phi_{g'} \\
    \sum^G_{g'=1}\Sigma_s^{g'\rightarrow 2}\phi_{g'} \\
    \vdots \\
    \sum^G_{g'=1}\Sigma_s^{g'\rightarrow G}\phi_{g'}
  \end{bmatrix},
  \mathbb{S}_1 \vec{\bm{J}} \equiv
  \begin{bmatrix}
    \sum^G_{g'=1}\Sigma_{s,1}^{g'\rightarrow 1}\vec{J}_{g'} \\
    \sum^G_{g'=1}\Sigma_{s,1}^{g'\rightarrow 2}\vec{J}_{g'} \\
    \vdots \\
    \sum^G_{g'=1}\Sigma_{s,1}^{g'\rightarrow G}\vec{J}_{g'}
  \end{bmatrix},
  \bm{Q}_0 \equiv
  \begin{bmatrix}
    \int_\mathcal{S} Q_1 d\hat{\Omega} \\
    \int_\mathcal{S} Q_2 d\hat{\Omega} \\
    \vdots \\
    \int_\mathcal{S} Q_G d\hat{\Omega}
  \end{bmatrix}, \nonumber
\end{gather}
%
and $\partial\mathcal{D}_v$ represents sections of $\partial\mathcal{D}$ which are vacuum
boundaries. We assume that $\bm{Q}$ is an isotropic source. We define the following expression
%
\begin{gather}
  \vec{\bm{r}}_1 \equiv \nabla\cdot(\vec{\vec{\bm{E}}}\bm{\Phi}) + \mathbb{L}_2\vec{\bm{J}} -
  \mathbb{S}_1\vec{\bm{J}} \label{eq:1st-moment}
\end{gather}
%
to simplify the balance equation. We also introduce the diffusion operator
%
\begin{gather}
  \mathbb{D}\nabla\bm{\Phi} \equiv
  \begin{bmatrix}
    D_1\nabla\phi_1 \\
    D_2\nabla\phi_2 \\
    \vdots \\
    D_G\nabla\phi_G
  \end{bmatrix},
\end{gather}
%
and recognize that under the diffusion approximation, the vacuum boundary condition is given as
%
\begin{gather}
  -\mathbb{D}\nabla\bm{\Phi}\cdot\hat{n}_b = \frac{\bm{\Phi}}{2} \hspace{0.5cm} \forall\vec{r}\in\partial
  \mathcal{D}_v.
\end{gather}
%
We can rewrite the balance equation as
%
\begin{multline}
  \left(\nabla\bm{\Phi}^*, \mathbb{D}\nabla\bm{\Phi}\right)_\mathcal{D}
  + \left(\mathbb{L}_2\bm{\Phi}^*,\bm{\Phi}\right)_\mathcal{D}
  + \left(\bm{\Phi}^*,\frac{\bm{\Phi}}{2}\right)_{\partial\mathcal{D}_v} \\
  + \underbrace{\left(\nabla\bm{\Phi}^*,\bm{\tau}\vec{\bm{r}}_1-\vec{\bm{J}}-\mathcal{D}\bm{\Phi}\right)_\mathcal{D}}_\text{(1)}
  + \underbrace{\left(\bm{\Phi}^*,\bm{J}^\text{out}-\frac{\bm{\Phi}}{2}\right)_{\partial\mathcal{D}_v}}_\text{(2)} \\
  = \left(\bm{\Phi}^*,\mathbb{S}_0\bm{\Phi}\right)_\mathcal{D}
  + \left(\bm{\Phi}^*,\bm{Q}_0\right)_\mathcal{D}. \label{eq:modified-diff}
\end{multline}
%
Eq.\ \ref{eq:modified-diff} is a modified neutron diffusion equation with
transport corrections provided by the terms labeled (1) and (2). Term (1) represents a drift term
with the drift vector defined as
%
\begin{gather}
  \vec{\bm{D}} \equiv \frac{\bm{\tau}\vec{\bm{r}}_1-\vec{\bm{J}}-\mathbb{D}\nabla\bm{\Phi}}{\bm{\Phi}}
\end{gather}
%
and Term (2) represents a boundary correction term with the boundary coefficient vector defined as
%
\begin{gather}
  \bm{\gamma} \equiv \frac{\bm{J}^\text{out}}{\bm{\Phi}}-\frac{1}{2}\bm{I}_3
\end{gather}
where $\bm{I}_3$ is the 3-D identity matrix.
With the $S_N$ angular discretization scheme, the drift and boundary correction vector components
can be evaluated as
%
\begin{gather}
  \vec{D}_g = \frac{\sum^{N_d}_{d=1}w_d\left(\tau_g\hat{\Omega}_d\hat{\Omega}_d\cdot\nabla\Psi_{g,d}
  + \left(\tau_g\Sigma_{t,g}-1\right)\hat{\Omega}_d\Psi_{g,d}
  - \tau_g\sum^G_{g'=1}\Sigma^{g'\rightarrow g}_{s,1}\hat{\Omega}_d\Psi_{g',d}
  - D_g\nabla\Psi_{g,d}\right)}{\sum^{N_d}_{d=1}w_d\Psi_{g,d}}, \label{eq:drift} \\
  \gamma_g =
  \frac{\sum_{\hat{\Omega}_d\cdot\hat{n}_b > 0}w_d |\hat{\Omega}_d\cdot\hat{n}_b |
  \Psi_{g,d}}{\sum^{N_d}_{d=1}w_d\Psi_{g,d}}. \label{eq:bound-coef}
\end{gather}

We limited the diffusion coefficients $D_g$ to a maximum value of 2.5 cm to improve numerical
stability in near-void regions.
This limiter does not affect the converged solution because the diffusion term containing $D_g$ is
asymptotically canceled by the transport correction terms as the \gls{HOLO} iterations converge.

\subsection{Hybrid $S_N$-Diffusion Method} \label{sec:hybrid-algorithm}

Transport correction formulations are commonly used in
iterative acceleration schemes or other hybrid methods to couple high-order neutron transport
calculations with low-order neutron diffusion calculations. However, the high-order calculations of
full reactor core models are computationally intensive and typically used for
time-independent neutronics simulations on \gls{HPC} clusters.

To reduce the computational cost of the high-order transport step, we propose reducing the
problem domain of the $S_N$ method to a \textit{correction subregion} containing the control rod
and its vicinity. Consequently, the hybrid $S_N$-diffusion method may retain accurate neutron flux
and current estimates around the control rod region from the $S_N$ method while making significant
computational cost savings by treating most of the reactor geometry with the neutron diffusion
method alone.

Hereafter, we refer to the $S_N$ system defined on the correction
region as the $S_N$ \textit{subsolver}. We define the full problem
domain and the correction subregion as $V_0$ and $V_1$, respectively, where
$V_1\subseteq V_0$.

%\subsection{$S_N$ Subsolver Boundary Conditions} \label{sec:sn-bc}

The $S_N$ subsolver is set up as a steady-state source problem with fission sources scaled by
the latest $k_\text{eff}$ estimate from the $k$-eigenvalue diffusion problem.
It also requires boundary sources along $\partial V_1$ derived from the
diffusion solution. We can obtain estimates for boundary fluxes
using the $P_1$ approximation to evaluate the neutron angular flux along
the discrete ordinates $\hat{\Omega}_d$ of the $S_N$ angular quadrature set as follows
%
\begin{align}
  \Psi_g(\vec{r},\hat{\Omega}_d) &\approx \frac{1}{4\pi}\left[\phi^\text{diff}_g(\vec{r})+
  3\hat{\Omega}_d\cdot\vec{J}^\text{diff}_g(\vec{r})\right] \nonumber \\
  &=\frac{1}{4\pi}\left[\phi^\text{diff}_g(\vec{r})-
  3\hat{\Omega}_d\cdot D_g\nabla\phi^\text{diff}_g(\vec{r})\right]
\end{align}
%
where the diff superscript denotes flux estimates from the latest neutron diffusion calculation.
With this relation, we can formulate the boundary source term for the $S_N$ subsolver by setting
the incident boundary flux source term,
$\Psi^\text{inc}_{g,d}$, in Eq.\ \ref{eq:boundary-source} to
%
\begin{gather}
  \Psi^\text{inc}_{g,d}(\vec{r}) \approx \frac{1}{w}
  \left[\phi^\text{diff}_g(\vec{r})-3\hat{\Omega}_d\cdot D_g\nabla\phi^\text{diff}_g(\vec{r})\right].
  \label{eq:p1-boundary}
\end{gather}

The $P_1$ approximation affects the accuracy of drift correction coefficients relative to reference
transport solutions near $\partial V_1$. Thus we propose setting a \textit{buffer region} in $V_1$
near $\partial V_1$, which can be interpreted as an extrapolation distance within which the drift
terms are discarded and not applied to the low-order neutron diffusion model. The buffer region
extends from $\partial V_1$ to a point in $V_1$ where drift components
%
\begin{equation}
  D_{g,j}=0\ \left(j=x,y,\mbox{or }z\right) \label{eq:cutoff}
\end{equation}
%
evaluated along the corresponding Cartesian axes. This preserves flux gradient smoothness across
the interface between corrected and uncorrected flux solutions.

Determining where the inner boundary of the buffer region is, based on Eq.\
\ref{eq:cutoff}, requires a precalculation step of running the hybrid $S_N$-diffusion method without
the buffer region to obtain drift correction distributions. Under-converged solutions obtained
after approximately two outer iterations should suffice for this step to reduce computational
expense. The buffer region can then be enforced by defining a small subregion $V_2$ in $V_1$,
containing the point where Eq.\ \ref{eq:cutoff} occurs and $D_{g,j}$ is continuous and monotonic
along $j$. These definitions enforce that the point where $D_{g,j}=0$, i.e., a $j$-intercept, is
unique in $V_2$. An adaptive buffer region boundary algorithm determines whether drift coefficients
should be applied to the low-order neutron diffusion model in $V_2$ using local drift and drift
gradient values as follows
%
\begin{gather}
  \left.D_{g,j}\right|_{V_2} =
  \begin{cases}
    D_{g,j} & \mbox{if } \left(j_1-j_0\right)D_{g,j}\frac{d D_{g,j}}{dj} < 0 \\
    0 & \mbox{if } \left(j_1-j_0\right)D_{g,j}\frac{d D_{g,j}}{dj} > 0
  \end{cases}, \label{eq:cutoff2}
  \shortintertext{where}
  \begin{align*}
    j_0 &= \mbox{ $j$-component of the reference center of }V_1, \\
    j_1 &= \mbox{ $j$-component of the local point in }V_2.
  \end{align*}
\end{gather}
%
Eq.\ \ref{eq:cutoff2} determines whether $D_{g,j}$ is approaching or past zero based on the
signedness of the local drift, drift gradient, and location in $V_2$ relative to the reference
center of $V_1$. The region between $V_2$ and $\partial V_1$ are automatically considered to be
within the buffer region.

With this buffer region implementation, $V_1$ is fixed while the inner buffer boundary shifts.
While more involved, this treatment allows for adaptive inner boundaries
in response to temperature reactivity feedback in non-isothermal multiphysics simulations. As for
$\partial V_1$, a practical guideline would be to place it about one neutron mean free path away
from the control rod to ensure the $P_1$ boundary conditions do not corrupt the transport solution
in the immediate vicinity of the control rod.

In our investigations, the buffer region also provided a two- to three-fold reduction in the total
number of outer iterations because the $S_N$ flux shape outside the buffer region
converges within the first two or three iterations. Since $\vec{D}_g$ is normalized by the flux
(Eq.\ \ref{eq:drift}), flux magnitude updates after the first few iterations do not impact it. In a
related context, Section \ref{sec:conv-tol} presents a sensitivity analysis of $k_\text{eff}$ with
respect to the $S_N$ subsolver convergence tolerance.

The iteration algorithm for the hybrid method is as follows:
%
\begin{enumerate}
  \item Start with an initial neutron diffusion $k$-eigenvalue calculation in $V_0$ using the
    neutron diffusion method without transport corrections.
  \item Pass the neutron diffusion flux solution along $\partial V_1$ to the
    $S_N$ subsolver and evaluate $S_N$ boundary source distribution using Eqs.\
    \ref{eq:boundary-source} and \ref{eq:p1-boundary}. Pass $k_\text{eff}$ for $S_N$ fission
    source scaling.
  \item Perform convergence check. If converged, terminate the iterations and output the latest
    flux and $k_\text{eff}$ from the neutron diffusion solver.
  \item Perform $S_N$ calculation and evaluate transport correction terms
    in $V_1$ using Eq.\ \ref{eq:drift}.
  \item Pass the transport correction terms to the neutron diffusion solver.
  \item Perform neutron diffusion calculation in $V_0$ with transport corrections in $V_1$.
  \item Repeat Steps 2 through 6 until the convergence criteria is met in Step 3.
\end{enumerate}
%
Convergence is declared when the $S_N$ subsolver residual, evaluated using Eq.\ \ref{eq:weak-form}
after updating the boundary sources, falls below the prescribed tolerance. This
criterion essentially represents convergence of the $S_N$-diffusion system, since the $S_N$
transport solution remains converged even after incorporating the updated boundary data from the
diffusion solver.

\subsection{Numerical Implementation} \label{sec:numerical-implementation}

This section details the numerical implementation of the hybrid $S_N$-diffusion method with
drift correction terms in Moltres.

Moltres \cite{lindsay_introduction_2018} is an open-source multiphysics reactor simulation software
developed for \gls{MSR} analysis and built on the \gls{MOOSE} framework
\cite{giudicelli30MOOSEEnabling2024}. \gls{MOOSE} enables highly scalable finite-element modeling
in part through its dependencies, libMesh \cite{kirkLibMeshLibraryParallel2006a} and PETSc
\cite{satishPETScUsersManual2019}.
\gls{MOOSE} also provides a developer-friendly interface for physics code development, native
multiphysics coupling to any \gls{MOOSE}-based application, and data postprocessing tools.
All source code for the hybrid $S_N$-diffusion model is openly available on the Moltres
GitHub repository\footnote{\url{https://github.com/arfc/moltres/}}.

Moltres presently comes with a Python script for extracting neutron group
constant data from OpenMC \cite{boyd_multigroup_2019} or Serpent 2 \cite{leppanen_serpent_2014}
simulation output files into a JSON format file. We updated the script to extract
the macroscopic total cross sections $\Sigma_{t,g}$ and the higher-order macroscopic Legendre
scattering moments $\Sigma^{g'\rightarrow g}_{s,l}$ for the new $S_N$ neutron transport model. A
new \texttt{Material} class called \texttt{MoltresSNMaterial} creates material
property variables for each group constant for the $S_N$ model to read during a neutronics
simulation and can interpolate between group constant datasets to simulate temperature-dependent
reactivity changes.

A new \texttt{MoltresUtils} class contains ``utility'' functions required by the $S_N$ model.
These include functions to: retrieve directions and weights defined by a
level-symmetric quadrature set of order $N$, retrieve reflecting directions at
reflecting boundaries, and compute spherical harmonic flux moments for anisotropic
scattering. The direction and weight values for the level-symmetric quadrature sets were
pre-computed and hard-coded into Moltres to reduce computational expense at runtime.

The \gls{SAAF} $S_N$ method defines $G$ energy groups and $N_d$ discrete directions for a total
of $G\times N_d$ angular flux variables $\Psi_{g,d}$. The Moltres $S_N$ model groups the
$\Psi_{g,d}$ variables by energy group into $G$ array variables. An array variable is a
set of standard field variables to help with dealing with setting up simulations with high
dimensionality. Each standard variable of an array variable is referred to as a component in the
array variable. In this work, we opted for
2nd-order Lagrange shape functions to approximate the $\Psi_{g,d}$ variables in the finite-element
analysis for a higher mesh convergence rate than 1st-order shape functions.
All $S_N$ simulations in this work employ the \gls{PJFNK} solver \cite{knoll_jacobian-free_2004}
with HYPRE-BoomerAMG \cite{hypre_hypre_2022} for preconditioning.

Moltres initializes the drift correction variables $\vec{D}_g$ as auxiliary array variables
with three components corresponding to their $x$, $y$, and $z$ components. Auxiliary variables
are field variables not directly solved for by the \gls{PJFNK} solver.
The drift and boundary correction variables are computed during
post-calculation stages using $\Psi_{g,d}$.
New kernel classes, \texttt{GroupDrift} and \texttt{VacuumCorrectionBC}, add
drift and boundary correction terms to the transport-corrected neutron diffusion solver.
Moltres leverages \texttt{MultiApp} and \texttt{Transfers} systems
\cite{gaston_physics-based_2015} to implement Picard iterative coupling between the $S_N$ and
neutron diffusion solvers for the hybrid $S_N$-diffusion method.
